\chapter*{Preface: how to read this book}
\addcontentsline{toc}{chapter}{Preface: how to read this book}

\textbf{Written for:} you, the author of the TFG, who has seen the results but has not yet had time to understand every method and wants to be able to extend the work alone.

\section*{What this book is}
The repository holds the code, the notes and the results of the first experiment of the TFG. Those files are written for someone who already knows the project. This book is written for someone who is \emph{learning} the project. It tells the story in order, explains each idea with small numbers you can check by hand, shows where every piece lives in the code, says honestly what is weak, and ends with recipes for doing your own analyses.

It replaces the two earlier guides in this folder (\file{tfg-guide.pdf}, a map, and \file{methods-study-guide.pdf}, a methods primer). They are kept; their material is reused here and expanded.

\section*{What this book is not}
It is not a statistics textbook. When a method needs more depth than a page or two, the text says so and points to a source in Chapter~\ref{ch:sources}, with a note on which chapter or section to read and how long it takes.

\section*{Three ways to read it}
\begin{description}
  \item[Cover to cover (about 6--8 hours).] The order is the logical order: why $\to$ data $\to$ tools $\to$ methods $\to$ code $\to$ results $\to$ doubts $\to$ sources $\to$ your own work.
  \item[Before a meeting (about 1 hour).] Chapter~\ref{ch:purpose}, the summary boxes at the start of each method in Chapter~\ref{ch:methods}, the table in Section~\ref{sec:headline}, and Chapter~\ref{ch:critique}.
  \item[To build something new (about 2 hours).] Chapter~\ref{ch:pipeline} (how the pieces connect), then Chapter~\ref{ch:building}. Come back to Chapter~\ref{ch:methods} for any method you reuse.
\end{description}

\section*{Boxes and conventions}
\begin{idea}
A blue box gives the idea of a section in one sentence. If you only remember one thing, remember this.
\end{idea}
\begin{pitfall}
A red box names a typical misreading or a typical mistake.
\end{pitfall}
\begin{tryit}
A green box asks you to do something small: compute by hand, run a snippet, change a number. Learning happens here. Answers to the exercises are in Appendix~\ref{app:exercises}.
\end{tryit}

File names are written like \file{04\_h1\_comment\_level.py}. Paths are relative to the \file{tfg/} folder. Square brackets such as \src{Mantel \& Haenszel 1959} refer to Chapter~\ref{ch:sources}.

\section*{Status of the claims}
Parts of the first analysis were done with an AI assistant (see \file{notes/method/AI-USE.md}). Every decision in the repository is tagged \emph{me}, \emph{AI, reviewed} or \emph{AI, not yet reviewed}. A goal of this book is that none of the decisions remains in the last category because you could not explain it. Chapter~\ref{ch:critique} lists what a line-by-line code review found; none of it changes a headline number, but several points limit what you may claim.

\section*{A note on how much is settled}
The headline results are numbers produced by code that was reviewed. The \emph{interpretation} of some of them is still open: for instance, whether the ``long-lived users'' result describes a trait of the users or just a property of how the group was defined (Section~\ref{sec:longlived}). Where interpretation is open, the book says so.
